\documentclass[a4paper,11pt]{article}
\usepackage{exam-style-842-unit}
\examinehead{842 信息化概论 · 第 2 单元练习册}{练习一\quad 数据管理与数据库基础}
\begin{document}

\unitcover{842 信息化概论 · 单元练习册}{第 2 单元\quad 数据管理、数据库与信息系统}{练习一\quad 数据管理与数据库基础（共 100 分）}{约 70 分钟}{数据与数据管理；数据管理发展的四个阶段；数据库、DBMS 与 DBS；数据模型三要素与概念模型（E-R）；层次、网状、关系三种数据模型；关系模型与键；三类完整性约束；三级模式与两级映像；视图；事务与 ACID；SQL；结构化数据与非结构化数据。}

\begin{knowmap}
\textbf{数据管理四阶段：}人工管理（20 世纪 50 年代中期以前）$\rightarrow$ 文件系统（20 世纪 50 年代末至 60 年代中期）$\rightarrow$ 数据库系统（20 世纪 60 年代末以后）$\rightarrow$ 高级数据库技术（20 世纪 80 年代以后）\\[2pt]
\textbf{数据抽象三级：}现实世界 $\rightarrow$ 概念模型（E-R：实体·属性·联系）$\rightarrow$ 逻辑模型（层次·网状·关系）$\rightarrow$ 物理模型\\[2pt]
\textbf{数据库三级模式：}外模式（用户视图）— 概念模式（全体数据的逻辑结构）— 内模式（物理存储结构）；两级映像：外模式/模式（逻辑独立性）、模式/内模式（物理独立性）\\[2pt]
\textbf{关系模型：}关系（二维表）· 元组（行）· 属性（列）· 域 · 键（候选键、主键、外键）；操作以关系代数为基础\\[2pt]
\textbf{完整性：}实体完整性（主键非空且唯一）· 参照完整性（外键）· 用户定义完整性\\[2pt]
\textbf{数据控制：}事务 ACID（原子性·一致性·隔离性·持久性）；视图是虚表，只存定义不存数据
\end{knowmap}

\begin{keyterms}
\textbf{数据}：描述客观事物属性及其相互关系的符号记录，有数值、文字、图形、声音等形式，是数据库中存储的基本对象。\\
\textbf{数据库（DB）}：以一定方式储存在一起、能为多个用户共享、具有尽可能小的冗余度、与应用程序彼此独立的数据集合。\\
\textbf{数据库管理系统（DBMS）}：位于用户与操作系统之间、用于科学地组织和存储数据、高效地获取和维护数据的大型系统软件。\\
\textbf{数据库系统（DBS）}：由数据库、数据库管理系统（及其开发工具）、应用程序和数据库管理员（DBA）组成的系统。\\
\textbf{数据管理}：对数据进行分类、组织、编码、存储、检索和维护的过程，其发展经历了人工管理、文件系统、数据库系统和高级数据库技术四个阶段。\\
\textbf{数据模型}：对现实世界数据特征的抽象，三要素是数据结构、数据操作和完整性约束。\\
\textbf{E-R 模型}：用实体、属性和联系描述概念模型的方法，实体用矩形、属性用椭圆、联系用菱形表示。\\
\textbf{关系模型}：用二维表（关系）表示实体及实体间联系的数据模型，对数据的操作通过关系代数实现，具有严格的数学基础。\\
\textbf{主键}：从候选键中选定的、能唯一标识关系中每个元组的属性或属性组，其取值必须唯一且不能为空。\\
\textbf{外键}：关系中的一个属性或属性组，其取值参照另一关系的主键，用于表示和约束实体间的联系。\\
\textbf{视图}：从一个或多个基本表（或其他视图）导出的虚表，数据库只存放视图的定义，不存放视图对应的数据。\\
\textbf{实体完整性}：关系中主键的值必须唯一，且不能取空值。\\
\textbf{参照完整性}：外键的取值必须是被参照关系中已存在的主键值，或者取空值。\\
\textbf{事务}：用户定义的一个不可分割的数据库操作序列，具有原子性、一致性、隔离性和持久性，即 ACID 特性。
\end{keyterms}

\begin{contrast}{文件系统 / 数据库系统（每行依次为：数据组织方式·冗余与共享·数据独立性）}
\citem{文件系统}{数据面向某一应用，若干数据文件彼此独立}{数据冗余度大、共享性差}{程序与数据之间独立性差，数据逻辑结构改变往往要修改程序}
\citem{数据库系统}{数据面向全局组织，可为多个用户共享}{冗余度尽可能小，易扩充、共享性高}{两级映像保证程序的逻辑独立性与物理独立性}
\end{contrast}

\begin{contrast}{层次 / 网状 / 关系模型（每行依次为：数据结构·能表示的联系·主要特点）}
\citem{层次模型}{树形结构，结点代表数据记录，连线表示从属关系}{一对多联系}{发展最早、技术较成熟，但结构呆板、查询须按路径导航}
\citem{网状模型}{有向图结构，结点间联系是任意的}{多对多联系}{比层次模型灵活，但结构复杂、使用不便}
\citem{关系模型}{二维表，用关系代数对数据进行操作}{通过公共属性表示实体间的联系}{结构简单、操作方便、有严格数学基础，20 世纪 80 年代后成为主流}
\end{contrast}

\begin{contrast}{外模式 / 概念模式 / 内模式（三级模式；每行依次为：面向对象·主要内容·独立性作用）}
\citem{外模式}{面向具体用户和应用}{用户能看到并使用的那部分数据的逻辑结构（子模式）}{通过外模式/模式映像实现逻辑独立性}
\citem{概念模式}{面向数据库整体}{全体数据的逻辑结构和特征}{是全部数据的公共视图，独立于存储结构}
\citem{内模式}{面向存储设备}{数据的物理存储结构和存取方法}{通过模式/内模式映像实现物理独立性}
\end{contrast}

\begin{contrast}{候选键 / 主键 / 外键（每行依次为：含义·取值约束·作用）}
\citem{候选键}{能唯一标识元组且不含多余属性的属性组}{可有多个，每个取值都唯一且非空}{作为选定主键的备选对象}
\citem{主键}{从候选键中选定的一个}{唯一且不能为空（实体完整性）}{唯一标识关系中的每一个元组}
\citem{外键}{参照另一关系主键的属性组}{取被参照关系的主键值或空值（参照完整性）}{表示并约束两个关系（实体）之间的联系}
\end{contrast}

\qtypename{一、单项选择题}{10 题，每题 3 分，共 30 分}
\q 数据管理的发展一般划分为四个阶段，其正确顺序是（\quad）
\choices{A. 文件系统$\rightarrow$人工管理$\rightarrow$数据库系统$\rightarrow$高级数据库技术}{B. 人工管理$\rightarrow$文件系统$\rightarrow$数据库系统$\rightarrow$高级数据库技术}
{C. 人工管理$\rightarrow$数据库系统$\rightarrow$文件系统$\rightarrow$高级数据库技术}{D. 人工管理$\rightarrow$文件系统$\rightarrow$高级数据库技术$\rightarrow$数据库系统}

\q “以一定方式储存在一起、能为多个用户共享、具有尽可能小的冗余度、与应用程序彼此独立的数据集合”指的是（\quad）
\choicesline{A. 数据}{B. 数据库}{C. 数据库管理系统}{D. 数据库系统}

\q 数据库系统的萌芽出现于（\quad）
\choicesline{A. 20 世纪 40 年代}{B. 20 世纪 50 年代}{C. 20 世纪 60 年代}{D. 20 世纪 80 年代}

\q 把数据按自然的层次关系组织起来、以反映数据之间隶属关系的数据模型是（\quad）
\choicesline{A. 层次模型}{B. 网状模型}{C. 关系模型}{D. 概念模型}

\q 对数据的操作通过关系代数实现、具有严格数学基础的数据模型是（\quad）
\choicesline{A. 层次模型}{B. 网状模型}{C. 关系模型}{D. 面向对象模型}

\q 关系数据库的标准语言 SQL 的全称是（\quad）
\choices{A. 结构化查询语言}{B. 标准化查询语言}{C. 系统查询语言}{D. 共享查询语言}

\q 结构化数据通常具有的两个重要特征是（\quad）
\choicesline{A. 排序和索引}{B. 加密和压缩}{C. 共享和传输}{D. 备份和恢复}

\q 数据库中只存放视图的定义，不存放视图对应的数据，因此视图被称为（\quad）
\choicesline{A. 基本表}{B. 虚表}{C. 索引}{D. 数据字典}

\q 在数据库三级模式中，描述数据物理存储结构和存取方法的是（\quad）
\choicesline{A. 外模式}{B. 概念模式}{C. 内模式}{D. 子模式}

\q 事务一旦提交，它对数据库的修改就是永久的，即使系统发生故障也不会丢失，这体现了事务的（\quad）
\choicesline{A. 原子性}{B. 一致性}{C. 隔离性}{D. 持久性}

\qtypename{二、多项选择题}{4 题，每题 4 分，共 16 分；每题有两个或两个以上正确项}
\q 关于数据库，下列说法正确的有（\quad）
\choices{A. 能为多个用户共享}{B. 具有尽可能小的冗余度}
{C. 与应用程序彼此独立}{D. 数据只能以文件形式存放}

\q 数据模型的三要素包括（\quad）
\choicesline{A. 数据结构}{B. 数据操作}{C. 完整性约束}{D. 存储结构}

\q 关于数据库的三级模式与两级映像，下列说法正确的有（\quad）
\choices{A. 外模式是用户能看到的数据视图}{B. 概念模式是全体数据的逻辑结构}
{C. 模式/内模式映像保证数据的物理独立性}{D. 外模式/模式映像保证数据的物理独立性}

\q 事务的 ACID 特性包括（\quad）
\choices{A. 原子性}{B. 一致性}{C. 隔离性}{D. 持久性}

\qtypename{三、判断题}{6 题，每题 2 分，共 12 分；正确的画“√”，错误的画“×”并写出改正后的表述}
\q 数据库中的数据冗余度越高越好，冗余越多数据越可靠。\quad（\quad）\quad 改正：\fixblank

\q 在关系模型中，主键的取值可以为空。\quad（\quad）\quad 改正：\fixblank

\q 视图是一个虚表，数据库中只存放视图的定义而不存放视图对应的数据。\quad（\quad）\quad 改正：\fixblank

\q 数据库三级模式结构中的外模式/模式映像保证了数据的物理独立性。\quad（\quad）\quad 改正：\fixblank

\q 数据库系统阶段的数据具有整体结构化特征，数据面向整个系统而不是某个应用。\quad（\quad）\quad 改正：\fixblank

\q 关系数据库对数据的操作通过关系代数实现，具有严格的数学基础。\quad（\quad）\quad 改正：\fixblank

\qtypename{四、配对题}{6 对，每对 1 分，共 6 分；把 A 栏与 B 栏对应内容配对}
\q 将下列概念与其对应内容配对，把 B 栏字母填入括号。
\par\vspace{0.2em}
\noindent\begin{tabular}{@{}p{0.42\textwidth}@{\hspace{1.2em}}p{0.46\textwidth}@{}}
（1）层次模型 \hfill（\quad） & A. 从基本表（或视图）导出的虚表 \\
（2）网状模型 \hfill（\quad） & B. 树形结构，结点间为一对多联系 \\
（3）关系模型 \hfill（\quad） & C. 有向图结构，结点间可表示多对多联系 \\
（4）视图 \hfill（\quad） & D. 二维表，用关系代数实现数据操作 \\
（5）主键 \hfill（\quad） & E. 参照另一关系主键的属性组 \\
（6）外键 \hfill（\quad） & F. 唯一标识元组且不能为空 \\
\end{tabular}\par\vspace{0.3em}

\qtypename{五、简答题}{3 题，每题 8 分，共 24 分}
\q 简述数据管理发展所经历的四个阶段及各阶段的主要特点。\anslines{2.4cm}

\q 简述数据库系统的三级模式结构，并说明两级映像如何保证数据的逻辑独立性和物理独立性。\anslines{2.4cm}

\q 什么是视图？视图有哪些优点？请至少写出四点。\anslines{2.4cm}

\newpage
\qtypename{六、综合题}{1 题，共 12 分}
\q 某县农业农村局要建立“种植户—地块—作物”数据库：一个种植户可以承包多个地块，一个地块可以种植多种作物，一种作物也可种在多个地块上；种植户有户编号、姓名、村组；地块有地块编号、面积、土壤类型；作物有作物编号、名称、生育期。
请回答：
（1）用 E-R 模型描述上述三个实体及其联系，指出联系的类型；
（2）把 E-R 模型转换为关系模式，并指出每个关系的主键和外键；
（3）说明应如何用实体完整性、参照完整性和用户定义完整性来保证该数据库数据的正确性；
（4）说明为该数据库建立视图有什么作用。\anslines{5cm}

\end{document}
